\chapter{Menus and prompts}\label{ch:menus}

ProPrep is driven entirely by menus and questions. This chapter shows you how to
read the main menu, how to answer the different kinds of questions, how to move
back and forth, and what the settings screens at the bottom of the menu do. It
ends with a short guide to Workflow Mode, the guided alternative to the Full
Menu.

\section{Reading the main menu}\label{sec:menus-main}

The Full Menu lists every tool, grouped in numbered sections that follow the
order of a typical preparation. Here is the start of it, just after starting a
new project:

\begin{screen}
Navigation: Main Menu


══ 1. STRUCTURE INPUT ══
  1a. ✓ Structure Loader
        Load protein structures from PDB, AlphaFold, AlphaFill, or local files
  1b. ○ Biological Assembly Generator
        Generate biological assembly from crystal asymmetric unit
        ⚠ Needs a loaded structure

══ 2. REDOX SITE DETECTION ══
  2a. ✓ Redox Site Detector
        Detect redox-active groups in protein structures
...
══ WORKSPACE & SETTINGS ══
  w. Workspace Options
  p. Preferences
  h. Help & Workflows
  f. Feedback & Support
  g. Switch to Workflow Mode
  x. Exit

Enter your choice (1a): §ans«1a»
\end{screen}

\begin{itemize}
  \item \textbf{Sections.} \textsf{1.~STRUCTURE INPUT}, \textsf{2.~REDOX SITE
  DETECTION}, \textsf{3.~STRUCTURE COMPARISON}, \textsf{4.~STRUCTURE
  PREPARATION}, \textsf{5.~SIMULATION SETUP \& EXECUTION}, \textsf{UTILITIES}
  and \textsf{WORKSPACE \& SETTINGS}.
  \item \textbf{Codes.} Each tool has a code made of the section number and a
  letter, such as \key{1a} or \key{4e}. Type the code and press \key{Enter} to
  open the tool. The settings at the bottom have one-letter codes. Appendix~\ref{app:menumap}
  lists them all.
  \item \textbf{The \textsf{Navigation} line} at the top of each screen is a
  \term{breadcrumb}: it shows where you are, for example
  \texttt{Navigation: Main Menu > Structure Loader}.
  \item \textbf{The default.} \prompt{Enter your choice (1a)} means that
  pressing \key{Enter} alone opens \key{1a}, the Structure Loader.
\end{itemize}

\subsection*{What the marks mean}

\begin{tabularx}{\linewidth}{@{}lL@{}}
\toprule
\textbf{Mark} & \textbf{Meaning} \\
\midrule
\texttt{✓} (green) & In the main menu: the tool \emph{can be used now}. It
does \emph{not} mean the step is done. \\
\texttt{○} (yellow) & The tool is not ready yet, because something it needs is
missing. \\
\texttt{⚠} & The line under a \texttt{○} tool that says what is missing, for
example \texttt{⚠ Needs a loaded structure} or
\texttt{⚠ Needs a topology — run Topology Generator first}. \\
\texttt{●} & Only \emph{inside} a tool's own menu: this action has been done
in this session. \\
\bottomrule
\end{tabularx}

\medskip
With nothing loaded, only \key{1a}, \key{2a}, \key{2b}, \key{3a} and \key{4g}
show \texttt{✓}. After you load a structure such as \pdb{1UBQ}, nearly every
tool turns \texttt{✓}. The tools that stay \texttt{○} wait for an earlier step:
\key{4d} and \key{5c} need detected redox sites, and \key{5b} needs a topology.

\begin{watchout}
The main menu never shows what you have already done. After filtering,
protonating and orienting a structure, \key{4a}, \key{4e} and \key{4f} still
show a plain \texttt{✓}. To see whether a step is done, open the tool and look
for \texttt{●} next to the action, as here after filtering:
\texttt{1.~● Filter PDB structure}.
\end{watchout}

You can open a \texttt{○} tool anyway. Its menu appears so you can look at its
options, but each blocked action shows \texttt{⚠} and does nothing when chosen.
Nothing breaks.

\begin{tip}
The Full Menu can also be drawn as two columns of panels. Start with
\typed{proprep --menu-grid}, or change it for good in Preferences
(Section~\ref{sec:menus-prefs}). It needs a wide window. The codes are the same.
\end{tip}

\section{Answering questions}\label{sec:menus-prompts}

Every question ends with a colon, and ProPrep waits for you to type an answer
and press \key{Enter}. Questions come in a few kinds.

\begin{tabularx}{\linewidth}{@{}p{0.34\linewidth}L@{}}
\toprule
\textbf{Looks like} & \textbf{What to type} \\
\midrule
\prompt{Enter your choice (1)} & A menu number or code. The value in
parentheses is the \term{default}: \key{Enter} alone accepts it. \\
\prompt{Remove hydrogen atoms? [y/n] (y)} & \key{y} for yes or \key{n} for no.
\key{Enter} accepts the default in parentheses. \\
\prompt{Select menu mode [1/2]} & One of the listed answers. With no
parentheses there is no default, so you must type an answer. \\
\prompt{Select chains (comma-separated indices, or 'all') (all)} & Several
numbers separated by commas, such as \typed{1,2}, or a word the question
offers. \\
\prompt{Enter PDB ID (e.g., 1ABC, 7XYZ)} & Free text: a code, a name, a number
or a file name. \\
\bottomrule
\end{tabularx}

\medskip
If you type something ProPrep does not expect, it says so and asks again:

\begin{screen}
Enter your choice (1a): §ans«9z»
Invalid choice: 9z. Please try again.
Hint: Use format like '1a', '3c', 'w', 'h', or 'x'
\end{screen}

Inside tools, a wrong answer gives a message such as:

\begin{screen}
Received 'zz', which is not one of the available options
\end{screen}

A yes/no question answers \texttt{Received '?'; please enter y or n}.

\begin{watchout}
The Full Menu accepts \textbf{lowercase only}. \typed{W} or \typed{4E} is
rejected as an invalid choice; type \key{w} or \key{4e}. (Workflow Mode accepts
both.) Also, a few questions that show no list of answers do not ask again
after a wrong answer: they leave the tool. If that happens, open the tool again.
\end{watchout}

\begin{watchout}
Always read the default before pressing \key{Enter}. Most defaults are the safe
choice, but not all. For example, pressing \key{Enter} at
\prompt{Submit this feedback to GitHub Issues? [y/n] (y)} sends your feedback,
and pressing \key{Enter} when picking a session to resume cancels the resume.
\end{watchout}

\section{Moving around: \texttt{b}, \texttt{m} and \texttt{x}}\label{sec:menus-nav}

Every tool menu ends with the same line:

\begin{screen}
[b]ack · [m]main · [x]exit
\end{screen}

\begin{tabularx}{\linewidth}{@{}lL@{}}
\toprule
\textbf{Key} & \textbf{What it does} \\
\midrule
\key{b} & Back one level. From a tool's first screen, this is the main menu. \\
\key{m} & Straight back to the main menu. \\
\key{x} & Exit ProPrep. It asks \prompt{Are you sure you want to exit? [y/n] (n)};
answering \key{n} takes you to the main menu, not back into the tool. \\
\bottomrule
\end{tabularx}

\medskip
Main-menu codes do not work inside a tool. Typing \typed{1a} inside the PDB
Filter gives \texttt{Received '1a', which is not one of the available options}.
Press \key{m} first, then type the code. Help pages end with
\prompt{Press Enter to return to main menu}, and the Preferences and Workspace
Options screens have a numbered ``Back to main menu'' item.

\section{Going back with \texttt{undo}}\label{sec:menus-undo}

If you gave a wrong answer, you do not need to start again. Type \typed{undo} at
any question. ProPrep lists your recent answers:

\begin{screen}
Enter your choice (1a): §ans«undo»

Recent answers in this session
# ┃ module                        ┃ question                     ┃ answer
━━╇━━━━━━━━━━━━━━━━━━━━━━━━━━━━━━━╇━━━━━━━━━━━━━━━━━━━━━━━━━━━━━━╇━━━━━━━━━━━━━━━━━━━━
0 │ Main Menu                     │ Select menu option           │ 1b  Biological ...
1 │ Biological Assembly Gener...  │ Select Biological Assemb...  │ 1  Generate biol...
2 │ Biological Assembly Gener...  │ Select Biological Assemb...  │ b  Back
3 │ Main Menu                     │ Select menu option           │ 9z
4 │ Main Menu                     │ Select menu option           │ W
Rewind to which answer? number to rewind to, [m]ore, or [x] cancel (4): §ans«1»
[r]e-answer it live, or [c]hange it and replay what followed? ([x] cancel) (r): §ans«r»
\end{screen}

Type the number of the answer you want to change, then press \key{Enter} to
accept \key{r} (re-answer it live). ProPrep quickly replays everything before
that answer and then asks that question again. A backup copy of the session
file is kept. Type \key{x} instead to cancel.

\begin{itemize}
  \item \key{undo} needs session recording. If you started ProPrep with
  \typed{--no-session}, ProPrep rejects it as an invalid choice.
  \item It cannot reach the menu-mode and project-name questions, which come
  before recording starts.
  \item Mistyped menu answers (like \texttt{9z} and \texttt{W} above) are recorded too, so
  they appear in the list.
\end{itemize}

\begin{watchout}
In version 1.23.0, the \key{c} choice (change the answer and replay what
followed) can be ignored for menu-type questions: ProPrep replays the old
answer. Use \key{r} and type the new answer yourself, or use the session editor.
Chapter~\ref{ch:sessions} explains \key{undo} and the session editor in full.
\end{watchout}

\section{Help \& Workflows (\texttt{h})}\label{sec:menus-help}

\key{h} shows one page called ``Common Workflows'': five typical sequences of
tools and a one-line description of each menu section. \key{Enter} returns to
the menu.

\begin{screen}
                    Common Workflows
...
  BASIC PROTEIN MD SIMULATION
   1a.Load → 3a.Filter → 3c.Fix → 3d.Protonate → 5b.tLEaP → 5c.MD
...
Press Enter to return to main menu...
\end{screen}

\begin{watchout}
In version 1.23.0, the codes on this help page are out of date. Typing them
opens the wrong tool: \key{3a} opens the Homology Searcher, not the PDB Filter.
Use the current codes:

\smallskip
\begin{tabularx}{\linewidth}{@{}lll@{}}
\toprule
\textbf{Tool} & \textbf{Help page says} & \textbf{Correct code} \\
\midrule
PDB Filter & 3a & \key{4a} \\
Structure Fixer & 3c & \key{4c} \\
Protonation State Analyzer & 3d & \key{4e} \\
Redox Site Preparer & 4a & \key{4d} \\
Force Field Parameterizer & 5a & \key{2c} \\
Topology Generator (tLEaP) & 5b & \key{5a} \\
Molecular Dynamics Manager & 5c & \key{5b} \\
QM/MM Preparator & 5d & \key{5c} \\
Structure Aligner & 2b & \key{3b} \\
\bottomrule
\end{tabularx}
\end{watchout}

Chapter~\ref{ch:bigpicture} gives an up-to-date overview of which tools to use
for which kind of project.

\section{Preferences (\texttt{p})}\label{sec:menus-prefs}

Preferences are settings that ProPrep remembers between runs. They are stored
in the file \file{~/.proprep/settings.json} in your home folder.

\begin{screen}
1. Change default menu mode (currently: Workflow Mode (guided))
2. Change full-menu layout (currently: List (single column))
3. Configure GitHub feedback (currently: not configured)
4. Check for new ProPrep releases at startup (currently: on)
5. Reset all preferences to defaults
6. Back to main menu
Enter your choice (1):
\end{screen}

\begin{tabularx}{\linewidth}{@{}p{0.27\linewidth}L@{}}
\toprule
\textbf{Option} & \textbf{What it does} \\
\midrule
1 Default menu mode & Choose \key{1} Workflow Mode, \key{2} Full Menu Mode or
\key{3} Ask each time. The default answer is always \key{1}, so type your
choice. It takes effect the next time you start ProPrep
(\texttt{Restart ProPrep for the change to take effect}). \\
2 Full-menu layout & \key{1} List (one tool per row, the default) or \key{2}
Grid (two columns of panels). Takes effect at once. \\
3 GitHub feedback & Stores a GitHub repository and access token so that
Feedback \& Support can post your reports (Chapter~\ref{ch:help}). \\
4 Release check & Turn the daily check for a newer ProPrep version
\key{1} On or \key{2} Off. \\
5 Reset & Puts every preference back to its default (asks first; the default
answer is \key{n}). This also removes a stored GitHub token. \\
6 Back & Returns to the main menu. \\
\bottomrule
\end{tabularx}

\section{Workspace Options (\texttt{w})}\label{sec:menus-workspace}

The \term{workspace} is what ProPrep holds in its memory during a run: the
structures you loaded, the files it made, and the results of each tool. The
Workspace Options screen lets you look at it and offers some technical tools.
Most users never need it.

\begin{tabularx}{\linewidth}{@{}p{0.34\linewidth}L@{}}
\toprule
\textbf{Option} & \textbf{What it does in 1.23.0} \\
\midrule
1 Show workspace status & A table of what is in memory. The names, such as
\texttt{rcsb\_pdb\_file}, are internal labels. \\
2 Show detailed workspace & Shows exactly the same table as option 1. \\
3 Save workspace & Writes \file{workspace.json}, but see the warning below. \\
4 Load workspace & Reads that file back; it restores almost nothing. \\
5 Reset workspace & Makes ProPrep forget everything loaded in this run (asks
first). Files on disk are kept. \\
6 Toggle debug mode & Shows extra technical messages; the same as starting with
\typed{--debug} (Chapter~\ref{ch:help}). \\
7 Show command history & The last ten menu-level actions, with times. \\
8--11 Command history & Save, load, replay or clear the command history. \\
12 Back to main menu & \\
\bottomrule
\end{tabularx}

\begin{watchout}
\textsf{Save workspace} does not save your work. In version 1.23.0 the file it
writes contains only \texttt{\{"pdb\_file": null\}}, even with a structure
loaded, and \textsf{Load workspace} brings nothing back. Your real record is the
session file, which ProPrep writes automatically (Chapter~\ref{ch:sessions}),
together with the files in your project folder (Chapter~\ref{ch:projectfolder}).
\end{watchout}

\begin{watchout}
\textsf{Reset workspace} also forgets the menu mode and the project folder. After
a reset, the breadcrumb shows \texttt{Navigation: Home}, and if you were in
Workflow Mode you land in the Full Menu. The reset is useful mainly when you want
ProPrep to stop using a structure from earlier in the run
(Section~\ref{sec:projectfolder-pickup}); starting a new project is usually
cleaner.
\end{watchout}

\section{Feedback \& Support (\texttt{f})}\label{sec:menus-feedback}

\key{f} opens a short form that turns a comment, question or problem report into
an issue on the developers' GitHub page. Chapter~\ref{ch:help} describes each
question and the GitHub setup it needs.

\begin{watchout}
At the end of the form, ProPrep asks
\prompt{Submit this feedback to GitHub Issues? [y/n] (y)}. The default is
\textbf{yes}: pressing \key{Enter} sends it. Type \key{n} if you do not want to
send. To leave the form early, press \key{Enter} at \prompt{Summary}; ProPrep
answers \texttt{Summary is required.} and \texttt{Feedback cancelled.}
\end{watchout}

\section{Workflow Mode}\label{sec:menus-workflow}

Workflow Mode is the guided version of the main menu. It shows only the tools
for one stage at a time, with a short suggestion, and you step from stage to
stage. You get it by choosing \key{1} at start-up, starting with
\typed{proprep --workflow}, or pressing \key{g} in the Full Menu.

\begin{screen}
╭───────────────────────────── WORKFLOW GUIDE ─────────────────────────────╮
│ Load:▶ → Detect:○ → Compare:○ → Fix:○ → Simulate:○                       │
│                                                                          │
│ Start by loading a structure from databases or local files and then      │
│ optionally applying symmetry transformations to generate a biological    │
│ assembly.                                                                │
╰──────────────────────────────────────────────────────────────────────────╯

══ 1. STRUCTURE INPUT ══
  1. ✓ Structure Loader
       Load protein structures from PDB, AlphaFold, AlphaFill, or local files
  2. ○ Biological Assembly Generator
       Generate biological assembly from crystal asymmetric unit
       ⚠ Needs a loaded structure

  v. Structure Viewer  |  u. Utilities...

[n]ext | [l]oad [d]etect [c]ompare [f]ix [s]imulate | [a]ll tools | e[x]it

Enter your choice (1):
\end{screen}

The guide bar at the top shows the five stages. \texttt{▶} marks the current
stage, \texttt{✓} a stage you have finished, \texttt{⊘} a stage you jumped over
and \texttt{○} a stage not reached yet.

\begin{tabularx}{\linewidth}{@{}lL@{}}
\toprule
\textbf{Stage} & \textbf{Tools (numbered within the stage)} \\
\midrule
Load & 1 Structure Loader, 2 Biological Assembly Generator \\
Detect & 1 Redox Site Detector, 2 Force Field Explorer, 3 Force Field
Parameterizer \\
Compare & 1 Homology Searcher, 2 Structure Aligner \\
Fix & 1 PDB Filter, 2 Amino Acid Mutator, 3 Structure Fixer, 4 Redox Site
Preparer, 5 Protonation State Analyzer, 6 Structure Orientator, 7 Membrane
Builder \\
Simulate & 1 Topology Generator, 2 Molecular Dynamics Manager, 3 QM/MM
Preparator, 4 Molecular Docking \\
\bottomrule
\end{tabularx}

\medskip
\begin{tabularx}{\linewidth}{@{}lL@{}}
\toprule
\textbf{Key} & \textbf{What it does} \\
\midrule
\key{1}, \key{2}, \ldots & Open a tool of the current stage. \\
\key{n} & Mark this stage done and go to the next one (not on Simulate). \\
\key{b} & Go back one stage (not on Load). \\
\key{l} \key{d} \key{c} \key{f} \key{s} & Jump to Load, Detect, Compare, Fix
or Simulate. \\
\key{v} & Structure Viewer (the Full Menu's \key{ua}). \\
\key{u} & Utilities: \key{w} Workspace Options, \key{p} Preferences, \key{h}
Help, \key{f} Feedback, \key{c} Cancel. \key{Enter} means Cancel here. \\
\key{a} & Switch to the Full Menu. \\
\key{x} & Exit (asks to confirm). \\
\bottomrule
\end{tabularx}

\medskip
After you load a structure and leave the Structure Loader with \key{b}, ProPrep
says \texttt{✓ Structure loaded! Advanced to Detect stage} and moves on by
itself.

\subsection*{Things to know about Workflow Mode}

\begin{itemize}
  \item The numbers restart at every stage: \key{1} is the Structure Loader on
  Load, but the PDB Filter on Fix and the Topology Generator on Simulate.
  \item \key{n} marks the stage you leave as done (\texttt{✓}) even if you did
  nothing there. Jumping forward with \key{s} marks the stages in between as
  skipped (\texttt{⊘}).
  \item Leaving the Structure Loader with \key{m} instead of \key{b} does not
  move you on to Detect. Press \key{n} yourself.
  \item After \key{b}, the stage you return to can show \texttt{✓} instead of
  \texttt{▶}. You are still on it.
  \item Switching to the Full Menu (\key{a}) and back (\key{g}) loses the stage
  you had reached.
  \item The Workflow help page (\key{u} then \key{h}) says that \key{s}
  ``skips'' a stage; in fact it jumps to Simulate. It also lists a ``Setup''
  stage that does not appear in the guide bar; topology building is on the
  Simulate stage.
\end{itemize}

None of these affect your files; they only affect the marks in the guide bar.
